\documentclass[10pt,a4paper]{amsart}
\usepackage{amsmath,amssymb,amsthm,graphicx}
% Retained minimal preamble for the Stacks Project chapter flat.tex (More on Flatness).
% Verbatim excerpt of the chapter preamble: Xy-pic setup (native preamble lines 14-18) and the native
% theorem declarations with the mathematical macros (native preamble lines 154-217). Copied mechanically;
% nothing added, removed or redefined. The wrapper class and the amsmath/amssymb/amsthm/graphicx base
% packages are supplied by the assembler, so the class/externaldocument part of the native preamble is omitted.
% For commutative diagrams we use Xy-pic
\usepackage[all]{xy}

% We use 2cell for 2-commutative diagrams.
\xyoption{2cell}
% Theorem environments.
%
\theoremstyle{plain}
\newtheorem{theorem}[subsection]{Theorem}
\newtheorem{proposition}[subsection]{Proposition}
\newtheorem{lemma}[subsection]{Lemma}

\theoremstyle{definition}
\newtheorem{definition}[subsection]{Definition}
\newtheorem{example}[subsection]{Example}
\newtheorem{exercise}[subsection]{Exercise}
\newtheorem{situation}[subsection]{Situation}

\theoremstyle{remark}
\newtheorem{remark}[subsection]{Remark}
\newtheorem{remarks}[subsection]{Remarks}

\numberwithin{equation}{subsection}

% Macros
%
\def\lim{\mathop{\mathrm{lim}}\nolimits}
\def\colim{\mathop{\mathrm{colim}}\nolimits}
\def\Spec{\mathop{\mathrm{Spec}}}
\def\Hom{\mathop{\mathrm{Hom}}\nolimits}
\def\Ext{\mathop{\mathrm{Ext}}\nolimits}
\def\SheafHom{\mathop{\mathcal{H}\!\mathit{om}}\nolimits}
\def\SheafExt{\mathop{\mathcal{E}\!\mathit{xt}}\nolimits}
\def\Sch{\mathit{Sch}}
\def\Mor{\mathop{\mathrm{Mor}}\nolimits}
\def\Ob{\mathop{\mathrm{Ob}}\nolimits}
\def\Sh{\mathop{\mathit{Sh}}\nolimits}
\def\NL{\mathop{N\!L}\nolimits}
\def\CH{\mathop{\mathrm{CH}}\nolimits}
\def\proetale{{pro\text{-}\acute{e}tale}}
\def\etale{{\acute{e}tale}}
\def\QCoh{\mathit{QCoh}}
\def\Ker{\mathop{\mathrm{Ker}}}
\def\Im{\mathop{\mathrm{Im}}}
\def\Coker{\mathop{\mathrm{Coker}}}
\def\Coim{\mathop{\mathrm{Coim}}}

% Boxtimes
%
\DeclareMathSymbol{\boxtimes}{\mathbin}{AMSa}{"02}

%
% Macros for moduli stacks/spaces
%
\def\QCohstack{\mathcal{QC}\!\mathit{oh}}
\def\Cohstack{\mathcal{C}\!\mathit{oh}}
\def\Spacesstack{\mathcal{S}\!\mathit{paces}}
\def\Quotfunctor{\mathrm{Quot}}
\def\Hilbfunctor{\mathrm{Hilb}}
\def\Curvesstack{\mathcal{C}\!\mathit{urves}}
\def\Polarizedstack{\mathcal{P}\!\mathit{olarized}}
\def\Complexesstack{\mathcal{C}\!\mathit{omplexes}}
% \Pic is the operator that assigns to X its picard group, usage \Pic(X)
% \Picardstack_{X/B} denotes the Picard stack of X over B
% \Picardfunctor_{X/B} denotes the Picard functor of X over B
\def\Pic{\mathop{\mathrm{Pic}}\nolimits}
\def\Picardstack{\mathcal{P}\!\mathit{ic}}
\def\Picardfunctor{\mathrm{Pic}}
\def\Deformationcategory{\mathcal{D}\!\mathit{ef}}

% --- appended 2026-10-12 (Chief): the project own definitions of the reference/slogan/history
% display environments, copied mechanically from the shared Stacks preamble lines 8-12
% (\usepackage{verbatim} plus the three \newenvironment lines, which the chapter preamble keeps
% verbatim from the project preamble). Nothing added, removed or redefined: these environments
% carry no mathematical content and the project itself typesets them as comments. Without these
% three lines a counted statement whose own text uses them fails with Environment slogan undefined.
\usepackage{verbatim}
\newenvironment{reference}{\comment}{\endcomment}
\newenvironment{slogan}{\comment}{\endcomment}
\newenvironment{history}{\comment}{\endcomment}
\usepackage{cleveref}
\crefname{originallemma}{original lemma}{original lemmas}
\Crefname{originallemma}{Original lemma}{Original lemmas}
\makeatletter
\newlabel{divisors-lemma-base-change-relative-assassin}{{05DC}{source}{}{sourcepointer-1}{}}
\newlabel{divisors-lemma-base-change-relative-assassin@cref}{{[originallemma][0][]05DC}{[1][0][]source}{}{}{}}
\makeatother
\crefname{originallemma}{original lemma}{original lemmas}
\Crefname{originallemma}{Original lemma}{Original lemmas}
\makeatletter
\newlabel{limits-lemma-descend-opens}{{01Z4}{source}{}{sourcepointer-2}{}}
\newlabel{limits-lemma-descend-opens@cref}{{[originallemma][0][]01Z4}{[1][0][]source}{}{}{}}
\makeatother
\crefname{originallemma}{original lemma}{original lemmas}
\Crefname{originallemma}{Original lemma}{Original lemmas}
\makeatletter
\newlabel{limits-lemma-inverse-limit-irreducibles}{{0CUG}{source}{}{sourcepointer-3}{}}
\newlabel{limits-lemma-inverse-limit-irreducibles@cref}{{[originallemma][0][]0CUG}{[1][0][]source}{}{}{}}
\makeatother
\crefname{originallemma}{original lemma}{original lemmas}
\Crefname{originallemma}{Original lemma}{Original lemmas}
\makeatletter
\newlabel{limits-lemma-limit-affine}{{01Z6}{source}{}{sourcepointer-4}{}}
\newlabel{limits-lemma-limit-affine@cref}{{[originallemma][0][]01Z6}{[1][0][]source}{}{}{}}
\makeatother
\crefname{originallemma}{original lemma}{original lemmas}
\Crefname{originallemma}{Original lemma}{Original lemmas}
\makeatletter
\newlabel{limits-lemma-limit-nonempty}{{01Z2}{source}{}{sourcepointer-5}{}}
\newlabel{limits-lemma-limit-nonempty@cref}{{[originallemma][0][]01Z2}{[1][0][]source}{}{}{}}
\makeatother
\crefname{originallemma}{original lemma}{original lemmas}
\Crefname{originallemma}{Original lemma}{Original lemmas}
\makeatletter
\newlabel{limits-lemma-separate}{{05BD}{source}{}{sourcepointer-6}{}}
\newlabel{limits-lemma-separate@cref}{{[originallemma][0][]05BD}{[1][0][]source}{}{}{}}
\makeatother
\crefname{originallemma}{original lemma}{original lemmas}
\Crefname{originallemma}{Original lemma}{Original lemmas}
\makeatletter
\newlabel{limits-proposition-characterize-locally-finite-presentation}{{01ZC}{source}{}{sourcepointer-7}{}}
\newlabel{limits-proposition-characterize-locally-finite-presentation@cref}{{[originallemma][0][]01ZC}{[1][0][]source}{}{}{}}
\makeatother
\crefname{originallemma}{original lemma}{original lemmas}
\Crefname{originallemma}{Original lemma}{Original lemmas}
\makeatletter
\newlabel{situation-pre-pure}{{05FW}{source}{}{sourcepointer-8}{}}
\newlabel{situation-pre-pure@cref}{{[originallemma][0][]05FW}{[1][0][]source}{}{}{}}
\makeatother
\title{An impurity of a module above a point descends along a directed limit of affine schemes}
\author{The Stacks Project Authors}
\date{stacks-project:flat.tex:local-snapshot-20261011; GNU FDL 1.2 or later}
\begin{document}
\maketitle
\noindent\textit{Editorial note.} All delivered mathematical text is the author's own, copied from the pinned Stacks Project chapter file flat.tex; no statement or proof step was rewritten, reordered or removed.\par\medskip
\noindent\textit{Editorial note.} Preamble provenance: the chapter's own retained minimal preamble (sources/stacks-flat/retained-preamble.tex, lines 1-85) is a verbatim excerpt of the native Xy-pic setup and of the native theorem declarations and mathematical macros, followed by the project's own reference/slogan/history display-environment definitions copied mechanically from the shared preamble; the wrapper class and the amsmath/amssymb/amsthm/graphicx base packages are supplied by the assembler, and extra\_packages is left empty because everything needed is already in the retained preamble. No macro was invented or redefined.\par\medskip
\noindent\textit{Editorial note.} Same-chapter pointers are located at their \textbackslash{}label lines inside the pinned chapter file; cross-chapter pointers are located in the project's own tag list staged in the snapshot (sources/stacks-chow/tags.txt), at the verbatim lines carrying their canonical tags.\par\medskip
\noindent\textit{Editorial note.} The counted proof span is the single proof environment at source lines 4582-4642, which the author internally organises as a 'First proof.' and a 'Second proof.' of the same assertion. Both are retained verbatim inside the span; nothing was dropped, and the environment is not a proof pointing at a later claim.\par\medskip
\noindent\textit{Allowed standard prerequisites (editorial).}\begin{itemize}
\item Control of the relative assassin under base change, and impurities. Divisors, label divisors-lemma-base-change-relative-assassin (proof line 4586), invoked to see that the image of the impurity point is again a relative associated point of the pulled back module; the situation defining impurities is situation-pre-pure (statement, proof line 4574). Reduces the task to showing only that the closure of the descended point misses the fibre over the descended specialisation, since the relative assassin condition is automatic.
\item Irreducible components of an inverse limit and emptiness of a limit of nonempty quasi-compact pieces. Limits, labels limits-lemma-inverse-limit-irreducibles (proof line 4596) and limits-lemma-limit-nonempty (proof line 4608). The first identifies the closure of the impurity point with the limit of the closures of its images, and the second turns the emptiness of the fibre product over the closed point into emptiness at some stage in the first proof.
\item Separation of a closed subset from a fibre and descent of open neighbourhoods and morphisms from a limit. Limits, labels limits-lemma-separate (proof line 4614), limits-lemma-descend-opens (proof line 4630), limits-lemma-limit-affine (proof line 4631) and limits-proposition-characterize-locally-finite-presentation (proof line 4635). The first produces the locally finitely presented replacement morphism and the closed subscheme disjoint from the special fibre, and the rest descend the affine open neighbourhood and the morphism from the replacement to a stage of the system in the second proof.
\end{itemize}\medskip
\ifdefined\hypertarget\hypertarget{sourcepointer-1}{}\fi
\noindent\textit{Source reference.} Original lemma 05DC: sources/stacks-chow/tags.txt, source lines 6950--6950. Navigation pointer for \textbackslash{}ref\{divisors-lemma-base-change-relative-assassin\} used inside the counted span (line 4586); the argument that consumes it is retained in the span.\par\medskip
\ifdefined\hypertarget\hypertarget{sourcepointer-2}{}\fi
\noindent\textit{Source reference.} Original lemma 01Z4: sources/stacks-chow/tags.txt, source lines 2579--2579. Navigation pointer for \textbackslash{}ref\{limits-lemma-descend-opens\} used inside the counted span (line 4630); the argument that consumes it is retained in the span.\par\medskip
\ifdefined\hypertarget\hypertarget{sourcepointer-3}{}\fi
\noindent\textit{Source reference.} Original lemma 0CUG: sources/stacks-chow/tags.txt, source lines 16101--16101. Navigation pointer for \textbackslash{}ref\{limits-lemma-inverse-limit-irreducibles\} used inside the counted span (line 4596); the argument that consumes it is retained in the span.\par\medskip
\ifdefined\hypertarget\hypertarget{sourcepointer-4}{}\fi
\noindent\textit{Source reference.} Original lemma 01Z6: sources/stacks-chow/tags.txt, source lines 2581--2581. Navigation pointer for \textbackslash{}ref\{limits-lemma-limit-affine\} used inside the counted span (line 4631); the argument that consumes it is retained in the span.\par\medskip
\ifdefined\hypertarget\hypertarget{sourcepointer-5}{}\fi
\noindent\textit{Source reference.} Original lemma 01Z2: sources/stacks-chow/tags.txt, source lines 2577--2577. Navigation pointer for \textbackslash{}ref\{limits-lemma-limit-nonempty\} used inside the counted span (line 4608); the argument that consumes it is retained in the span.\par\medskip
\ifdefined\hypertarget\hypertarget{sourcepointer-6}{}\fi
\noindent\textit{Source reference.} Original lemma 05BD: sources/stacks-chow/tags.txt, source lines 6881--6881. Navigation pointer for \textbackslash{}ref\{limits-lemma-separate\} used inside the counted span (line 4614); the argument that consumes it is retained in the span.\par\medskip
\ifdefined\hypertarget\hypertarget{sourcepointer-7}{}\fi
\noindent\textit{Source reference.} Original lemma 01ZC: sources/stacks-chow/tags.txt, source lines 2587--2587. Navigation pointer for \textbackslash{}ref\{limits-proposition-characterize-locally-finite-presentation\} used inside the counted span (line 4635); the argument that consumes it is retained in the span.\par\medskip
\ifdefined\hypertarget\hypertarget{sourcepointer-8}{}\fi
\noindent\textit{Source reference.} Original lemma 05FW: sources/stacks-flat/pinned-original-source.tex, source lines 4473--4473. Navigation pointer for \textbackslash{}ref\{situation-pre-pure\} used inside the counted span (line 4574); the argument that consumes it is retained in the span.\par\medskip
\begin{lemma}
\label{lemma-impure-limit}
In Situation \ref{situation-pre-pure}.
Let $(g : T \to S, t' \leadsto t, \xi)$ be an impurity of
$\mathcal{F}$ above $s$. Assume $T = \lim_{i \in I} T_i$
is a directed limit of affine schemes over $S$. Then for
some $i$ the triple $(T_i \to S, t'_i \leadsto t_i, \xi_i)$
is an impurity of $\mathcal{F}$ above $s$.
\end{lemma}

\begin{proof}
The notation in the statement means this: Let $p_i : T \to T_i$
be the projection morphisms, let $t_i = p_i(t)$ and $t'_i = p_i(t')$.
Finally $\xi_i \in X_{T_i}$ is the image of $\xi$. By
Divisors, Lemma \ref{divisors-lemma-base-change-relative-assassin}
it is true that $\xi_i$ is a point of the relative
assassin of $\mathcal{F}_{T_i}$ over $T_i$. Thus the only point is to
show that $\overline{\{\xi_i\}} \cap X_{t_i} = \emptyset$ for some $i$.

\medskip\noindent
First proof. Let $Z_i = \overline{\{\xi_i\}} \subset X_{T_i}$
and $Z = \overline{\{\xi\}} \subset X_T$
endowed with the reduced induced scheme structure.
Then $Z = \lim Z_i$ by
Limits, Lemma \ref{limits-lemma-inverse-limit-irreducibles}.
Choose a field $k$ and a morphism $\Spec(k) \to T$ whose image is $t$.
Then
$$
\emptyset =
Z \times_T \Spec(k) = (\lim Z_i) \times_{(\lim T_i)} \Spec(k)
= \lim Z_i \times_{T_i} \Spec(k)
$$
because limits commute with fibred products (limits commute with limits).
Each $Z_i \times_{T_i} \Spec(k)$ is quasi-compact because $X_{T_i} \to T_i$
is of finite type and hence $Z_i \to T_i$ is of finite type.
Hence $Z_i \times_{T_i} \Spec(k)$ is empty for some $i$ by
Limits, Lemma \ref{limits-lemma-limit-nonempty}.
Since the image of the composition $\Spec(k) \to T \to T_i$ is $t_i$
we obtain what we want.

\medskip\noindent
Second proof. Set $Z = \overline{\{\xi\}}$. Apply
Limits, Lemma \ref{limits-lemma-separate}
to this situation to obtain an open neighbourhood
$V \subset T$ of $t$, a commutative diagram
$$
\xymatrix{
V \ar[d] \ar[r]_a & T' \ar[d]^b \\
T \ar[r]^g & S,
}
$$
and a closed subscheme $Z' \subset X_{T'}$ such that
\begin{enumerate}
\item the morphism $b : T' \to S$ is locally of finite presentation,
\item we have $Z' \cap X_{a(t)} = \emptyset$, and
\item $Z \cap X_V$ maps into $Z'$ via the morphism $X_V \to X_{T'}$.
\end{enumerate}
We may assume $V$ is an affine open of $T$, hence by
Limits, Lemmas \ref{limits-lemma-descend-opens} and
\ref{limits-lemma-limit-affine}
we can find an $i$ and an affine open $V_i \subset T_i$ with
$V = f_i^{-1}(V_i)$. By
Limits,
Proposition \ref{limits-proposition-characterize-locally-finite-presentation}
after possibly increasing $i$ a bit we can find a morphism
$a_i : V_i \to T'$ such that $a = a_i \circ f_i|_V$.
The induced morphism $X_{V_i} \to X_{T'}$ maps $\xi_i$ into
$Z'$. As $Z' \cap X_{a(t)} = \emptyset$ we conclude that
$(T_i \to S, t'_i \leadsto t_i, \xi_i)$ is an impurity of
$\mathcal{F}$ above $s$.
\end{proof}

\end{document}
